% s-317-p317-m002; original PDF 317, 318
\Aphorism{s-317-p317-m002}{372}
Dass Spinoza sich der mathematischen Beweisform bediente, war gewiss überflüssig; aber noch viel überflüssiger, dass Kant die Unmöglichkeit dieser Form für philosophische Zwecke beweisen wollte, aus „reinen Begriffen“, wie man sich denken kann. Jede beliebige Schlusskette lässt sich, wenn man durchaus will, \Foreign{more geometrico} vertragen; es versteht sich von selbst, dass an Natürlichkeit und Eleganz dadurch nichts gewonnen, an Papier und Worten nichts gespart wird. Das wissen nicht nur Philosophen, sondern auch die Mathematiker selbst und beschränken die schwerfällige Handhabung von Axiomen, Lehrsätzen und Beweisen, als unentbehrliche Gymnastik für Anfänger, auf den Schulunterricht. Aber solche äusseren, darstellerischen Gründe hatte Kant nicht im Auge; wo wir an Spinoza einen Schönheitsfehler bemerken, hätte er am liebsten einen „transcendentalen Irrthum“ aufgestöbert, eine Principienfrage von ungemeiner Wichtigkeit, einen Cardinalpunkt, bei dem man sich um des ewigen Seelenheils willen nicht vergreifen dürfe! Wenn irgendworin, so hat Kant gerade in dieser masslosen Aufbauschung und Verallgemeinerung recht harmloser Dinge die Erbschaft der Religion angetreten; Beide reden nie anders als hyperbolisch und sehen die Mücke als Elephanten, Beiden fehlt der Blick für \textit{\Greek{ἀδιάφορα}}, für das weite breite sonnige Gebiet des Indifferenten.
\par

